\chapter{A rendszer tervezése}
\label{ch:tervezes}
% Terjedelem: kb. 7–8 oldal (~12 000 karakter)

\section{Követelmények}
\subsection{Funkcionális követelmények}
% Szerepkörök: hallgató, oktató (meghívókóddal regisztrál).
% Oktató: vizsga létrehozása (időablak + időtartam), feladatok típusonként, élő monitor,
%   eredmények, kézi pontozás + visszajelzés, CSV export, újraértékelés.
% Hallgató: vizsga indítása, visszaszámláló, automatikus mentés, próbafuttatás, beadás, eredmény.
\todo{Szöveg}

\subsection{Nem funkcionális követelmények}
% Biztonság (megoldókulcs sosem jut a klienshez, sandbox), használhatóság, hordozhatóság (Docker).
\todo{Szöveg}

\section{Architektúra}
% Háromrétegű felépítés: Next.js kliens – FastAPI + Socket.IO szerver – PostgreSQL,
% mellette a Docker sandbox. Ide jön egy ábra (abrak/architektura.png):
\todo{Szöveg, ábra: \ref{fig:architektura}.~ábra}

\begin{figure}[H]
  \centering
  % \includegraphics[width=0.9\textwidth]{architektura}
  \fbox{\parbox[c][4cm][c]{0.8\textwidth}{\centering Helyőrző: architektúra ábra}}
  \caption{A rendszer architektúrája}
  \label{fig:architektura}
\end{figure}

\section{Adatmodell}
% Entitások: User, Exam, Question, Attempt, Answer – ER-diagram + magyarázat.
% Fontos tervezési döntések: egy hallgató/vizsga = egy Attempt (UNIQUE), egy válasz/feladat (UNIQUE),
% answer_key JSON típusonként eltérő szerkezettel, manual_score felülírja az auto_score-t.
A rendszer adattábláit a \ref{tab:adatmodell}.~táblázat foglalja össze.

\begin{table}[H]
  \centering
  \caption{Az adatmodell táblái}
  \label{tab:adatmodell}
  \begin{tabularx}{\textwidth}{@{}lX@{}}
    \toprule
    Tábla & Szerep és fontosabb mezők \\
    \midrule
    \texttt{users}     & felhasználók; e-mail, név, jelszó-hash, szerepkör (\texttt{student}/\texttt{teacher}) \\
    \texttt{exams}     & vizsgák; cím, időtartam (perc), kezdési és záró időpont, létrehozó oktató \\
    \texttt{questions} & feladatok; típus, pontszám, válaszlehetőségek, megoldókulcs (JSON) \\
    \texttt{attempts}  & kitöltések; kezdés, határidő, befejezés, szabálysértések száma \\
    \texttt{answers}   & válaszok; válasz (JSON), automatikus és kézi pontszám, visszajelzés, állapot \\
    \bottomrule
  \end{tabularx}
\end{table}

\section{A REST API terve}
\begin{table}[H]
  \centering
  \caption{A REST API fontosabb végpontjai}
  \label{tab:api}
  \small
  \begin{tabularx}{\textwidth}{@{}lp{6.2cm}X@{}}
    \toprule
    Metódus & Végpont & Funkció \\
    \midrule
    POST   & \texttt{/auth/register}, \texttt{/auth/login} & regisztráció, bejelentkezés (JWT) \\
    POST   & \texttt{/exams/} & vizsga létrehozása (oktató) \\
    POST   & \texttt{/exams/\{id\}/questions} & feladat hozzáadása \\
    POST   & \texttt{/exams/\{id\}/attempt} & kitöltés indítása (hallgató) \\
    PUT    & \texttt{/attempts/\{id\}/answers/\{qid\}} & válasz mentése \\
    POST   & \texttt{/attempts/\{id\}/run} & kód próbafuttatása \\
    POST   & \texttt{/attempts/\{id\}/finish} & végleges beadás \\
    GET    & \texttt{/exams/\{id\}/results(.csv)} & eredmények, CSV export \\
    PATCH  & \texttt{/answers/\{id\}/grade} & kézi pontozás \\
    POST   & \texttt{/exams/\{id\}/regrade} & újraértékelés \\
    \bottomrule
  \end{tabularx}
\end{table}

\section{Biztonsági terv}
% Fenyegetések és ellenintézkedések: jogosultság (csak a saját vizsga), megoldókulcs elrejtése,
% Socket.IO JWT-hitelesítés, sandbox korlátok.
\todo{Szöveg}
